\documentclass[12pt]{article}

\input{~/School/Notes/header.tex}

\title{Домашнее задание по алгебре}
\author{Денисов Илья}
\date{}

\begin{document}
	\maketitle

	\section*{Задача 1}

	\noindent
	\textbf{Условие.} Многочлен $f$ степени $101$ таков, что
	\[
		f(n)=\lfloor 1{,}01n\rfloor
	\]
	для каждого целого $n$ от $99$ до $200$. Найдите $f(0)$.

	\begin{proof}[Решение]
		Сначала выясним, чему равны заданные значения. Поскольку $n$~--- целое число,
		\[
			\lfloor 1{,}01n\rfloor
			=\left\lfloor n+\frac{n}{100}\right\rfloor
			=n+\left\lfloor\frac{n}{100}\right\rfloor.
		\]
		Поэтому
		\[
			f(99)=99,\qquad
			f(n)=n+1\quad(100\leqslant n\leqslant199),\qquad
			f(200)=202.
		\]

		Рассмотрим разность $f(x)-(x+1)$. Она обращается в ноль при каждом целом
		$x$ от $100$ до $199$. Таких чисел ровно $100$. Обозначим
		\[
			W(x)=\prod_{j=100}^{199}(x-j).
		\]
		Следовательно, многочлен $f(x)-(x+1)$ делится на $W(x)$. Так как
		$\deg f=101$, а $\deg W=100$, частное имеет степень $1$. Значит,
		\[
			f(x)=x+1+W(x)(Ax+B)
		\]
		при некоторых числах $A$ и $B$.

		Найдём $A$ и $B$ из оставшихся двух значений: при $x=99$ и $x=200$.
		В обеих точках произведение $W$ легко вычислить:
		\begin{align*}
			W(99)
				&=(99-100)(99-101)\cdots(99-199)=100!,\\
			W(200)
				&=(200-100)(200-101)\cdots(200-199)=100!.
		\end{align*}
		В первом произведении $100$ отрицательных множителей, поэтому его
		знак положителен.

		Подставляя $x=99$, получаем
		\[
			99=100+100!(99A+B),
			\qquad\text{то есть}\qquad
			99A+B=-\frac{1}{100!}.
		\]
		При $x=200$ получаем
		\[
			202=201+100!(200A+B),
			\qquad\text{то есть}\qquad
			200A+B=\frac{1}{100!}.
		\]
		Вычтем первое равенство из второго:
		\[
			101A=\frac{2}{100!},
			\qquad
			A=\frac{2}{101\cdot100!}.
		\]
		Теперь из $200A+B=1/100!$ следует
		\[
			B=\frac{1}{100!}-\frac{400}{101\cdot100!}
			 =-\frac{299}{101\cdot100!}.
		\]
		Таким образом,
		\[
			f(x)=x+1+\frac{2x-299}{101\cdot100!}
				\prod_{j=100}^{199}(x-j).
		\]

		Осталось подставить $x=0$. Здесь снова будет $100$ отрицательных
		множителей, так что
		\[
			W(0)=\prod_{j=100}^{199}(-j)
				=100\cdot101\cdots199
				=\frac{199!}{99!}.
		\]
		Следовательно,
		\[
			f(0)
				=1-\frac{299}{101\cdot100!}\frac{199!}{99!}
				=1-\frac{299}{101}\binom{199}{100}.
		\]
	\end{proof}

	\noindent
	\textbf{Ответ:}
	\[
		\boxed{f(0)=1-\frac{299}{101}\binom{199}{100}}.
	\]

	%==================================================
	\section*{Задача 2}

	\noindent
	\textbf{Условие.} Дан многочлен $f$ третьей степени. Назовём различные
	числа $a$ и $b$ друзьями, если $f(a)=b$ и $f(b)=a$. Докажите, что есть
	не более четырёх пар друзей.

	\begin{proof}
		Если $a$ и $b$~--- друзья, то
		\[
			f(f(a))=f(b)=a,\qquad f(f(b))=f(a)=b.
		\]
		Иными словами, оба числа $a$ и $b$ являются корнями многочлена
		\[
			H(x)=f(f(x))-x.
		\]
		Однако при подсчёте всех корней $H$ нужно учитывать и числа,
		которые друзьями ни с кем не являются: неподвижные точки $f(x)=x$
		тоже удовлетворяют равенству $f(f(x))=x$.

		Чтобы отделить их от друзей, положим
		\[
			h(x)=f(x)-x.
		\]
		Заметим, что
		\[
			H(x)=f(f(x))-x
			 =\bigl(f(f(x))-f(x)\bigr)+\bigl(f(x)-x\bigr).
		\]
		Для любого многочлена $f$ разность $f(u)-f(v)$ делится на $u-v$.
		Применим это при $u=f(x)$ и $v=x$. Тогда
		$f(f(x))-f(x)$ делится на $f(x)-x=h(x)$. Второе слагаемое также
		равно $h(x)$. Значит, $H(x)$ делится на $h(x)$:
		\[
			H(x)=h(x)Q(x)
		\]
		для некоторого многочлена $Q$.

		Поскольку $\deg f=3$, имеем
		\[
			\deg H=9,\qquad \deg h=3,\qquad \deg Q=6.
		\]
		Теперь пусть $a$ и $b$~--- друзья. Тогда $H(a)=H(b)=0$. При этом
		они различны, поэтому
		\[
			h(a)=f(a)-a=b-a\ne0,\qquad
			h(b)=f(b)-b=a-b\ne0.
		\]
		Из равенства $H=hQ$ следует, что
		\[
			Q(a)=Q(b)=0.
		\]
		Итак, каждой паре друзей соответствуют \emph{два различных корня}
		многочлена $Q$.

		Разные пары друзей не могут содержать одно и то же число. Действительно,
		если $a$ дружит с $b$, то его друг однозначно определён:
		$b=f(a)$. Поэтому $m$ разных пар дают $2m$ разных корней $Q$.
		Но у многочлена шестой степени не более шести различных корней.
		Следовательно,
		\[
			2m\leqslant6,\qquad m\leqslant3.
		\]

		Мы доказали даже более сильное утверждение: пар друзей не более
		\textbf{трёх}. В частности, их не более четырёх, как требовалось.
	\end{proof}

	%==================================================
	\section*{Задача 3}

	\noindent
	\textbf{Условие.} Дана интерполяционная задача с $n$ комплексными
	узлами, мнимые части которых положительны. Докажите, что найдётся
	многочлен с вещественными коэффициентами, решающий эту задачу; его
	степень не обязательно должна быть меньше $n$.

	\begin{proof}
		Запишем задачу явно. Пусть заданы попарно различные узлы
		$z_1,\ldots,z_n$ и произвольные комплексные значения
		$w_1,\ldots,w_n$. Требуется найти многочлен $P$ с вещественными
		коэффициентами такой, что
		\[
			P(z_j)=w_j\qquad(j=1,\ldots,n).
		\]

		Воспользуемся условием $\operatorname{Im}z_j>0$. Сопряжённые числа
		$\overline{z_1},\ldots,\overline{z_n}$ лежат в нижней полуплоскости,
		поскольку
		\[
			\operatorname{Im}\overline{z_j}
				=-\operatorname{Im}z_j<0.
		\]
		Поэтому ни один из исходных узлов не совпадает с сопряжённым
		узлом. Кроме того, сопряжённые узлы попарно различны. Таким образом,
		все $2n$ чисел
		\[
			z_1,\ldots,z_n,\overline{z_1},\ldots,\overline{z_n}
		\]
		попарно различны.

		Добавим к исходным условиям ещё $n$ условий:
		\[
			P(\overline{z_j})=\overline{w_j}
			\qquad(j=1,\ldots,n).
		\]
		По теореме об интерполяции через $2n$ попарно различных узлов
		существует единственный многочлен $P$ с комплексными
		коэффициентами степени не выше $2n-1$, удовлетворяющий всем
		этим $2n$ условиям.

		Докажем, что его коэффициенты на самом деле вещественные.
		Если
		\[
			P(x)=c_0+c_1x+\dots+c_mx^m,
		\]
		рассмотрим многочлен
		\[
			P^*(x)=\overline{c_0}
				+\overline{c_1}x+\dots+\overline{c_m}x^m.
		\]
		Для любого комплексного $u$ выполняется
		\[
			P^*(u)=\overline{P(\overline{u})}.
		\]
		Поэтому в исходных узлах
		\[
			P^*(z_j)
				=\overline{P(\overline{z_j})}
				=\overline{\overline{w_j}}
				=w_j,
		\]
		а в сопряжённых узлах
		\[
			P^*(\overline{z_j})
				=\overline{P(z_j)}
				=\overline{w_j}.
		\]
		Итак, $P^*$ удовлетворяет тем же $2n$ интерполяционным
		условиям, что и $P$. В силу единственности $P^*=P$. Значит,
		$c_j=\overline{c_j}$ для каждого $j$, то есть все коэффициенты
		$P$ вещественны.

		Полученный многочлен решает исходную задачу и имеет степень
		не выше $2n-1$.
	\end{proof}

	%==================================================
	\section*{Задача 5}

	\noindent
	\textbf{Условие.} Рассмотрим интерполяционную задачу
	$f(i)=a_i$ в поле вычетов по модулю $p$, где
	$i=0,1,\ldots,p-1$. Проверьте, что интерполяционная формула упрощается до:
	\[
		f(x)=(x-x^p)\sum_{i=0}^{p-1}\frac{a_i}{x-i}.
	\]

	\begin{proof}[Проверка формулы]
		Пусть
		\[
			W(x)=\prod_{j=0}^{p-1}(x-j)\in\mathbb F_p[x].
		\]
		Все элементы $\mathbb F_p$ удовлетворяют равенству $i^p=i$.
		Поэтому каждый из них является корнем $x^p-x$. Многочлены
		$W(x)$ и $x^p-x$ имеют степень $p$, старший коэффициент $1$
		и одни и те же $p$ различных корней. Следовательно,
		\[
			W(x)=x^p-x.
		\]

		Формула Лагранжа для данных $f(i)=a_i$ имеет вид
		\[
			f(x)=
			\sum_{i=0}^{p-1}
				a_i
				\frac{\displaystyle\prod_{j\ne i}(x-j)}
				     {\displaystyle\prod_{j\ne i}(i-j)}.
		\]
		Числитель дроби равен $W(x)/(x-i)$. Вычислим её знаменатель.
		Так как
		\[
			W(x)=(x-i)\prod_{j\ne i}(x-j),
		\]
		при подстановке $x=i$ в производную получаем
		\[
			W'(i)=\prod_{j\ne i}(i-j).
		\]
		С другой стороны, в поле характеристики $p$
		\[
			W'(x)=(x^p-x)'=p x^{p-1}-1=-1.
		\]
		Значит, $\prod_{j\ne i}(i-j)=-1$ для каждого $i$.
		Подставляя это в формулу Лагранжа, получаем
		\[
			f(x)
				=-\sum_{i=0}^{p-1}a_i\frac{W(x)}{x-i}
				=-(x^p-x)\sum_{i=0}^{p-1}\frac{a_i}{x-i}.
		\]
		Это многочлен степени не выше $p-1$, как и должен быть
		интерполяционный многочлен для $p$ узлов.
	\end{proof}
\end{document}